\scanpage{17}
\mspage{8}
Si $\mathfrak{D}$ est un groupe isomorphe à $D_4$, l'ens $S$ des éléments d'ordre $2$ \deleted{\unsure{non centraux}}\inserted{$\ne\underline{a}$ (où $\underline{a}$ est} l'unique élément central $\ne1$) \deleted{\unsure{est}} a quatre éléments, \deleted{\unsure{qui permet déjà de voir par quelle}} \inserted{sur lequel on a une involution} sans point \deleted{\unsure{fixe}} fixe par
\[
x\longmapsto x\unsure{\cdot|+}\underline{a}\underset{\mathrm{def}}{=}x'
\]
donc à $\mathfrak{D}$ correspond un carré $Q_{\mathfrak{D}}$, dont l'ens des sommets est $S$, et l'\deleted{\unsure{involution de}}\inserted{action de} l'antipodisme sur $S$ est l'involution précédente.
Mais le foncteur de groupoïdes $\mathfrak{D}\mapsto Q_{\mathfrak{D}}$
\[
(\text{groupes isom. à }D_4)\longrightarrow(\text{carrés combinatoires})
\]
est une équivalence de groupoïdes, car on reconstitue $\mathfrak{D}$ à partir de $(S,\underline{a}_S:\inserted{x\mapsto x'})$ \deleted{\unsure{prenant}}\inserted{par} la formule
\[
\mathfrak{D}=\left\{
\underbrace{x\in S,\ \underline{a}}_{\text{générateurs}}\ \middle|\
\underbrace{[x,y]}_{\substack{\text{commutateur}\\xyx^{-1}y^{-1}}}
=\underline{a}\ \text{si }y\ne x,x'
\right\}.
\]
\margin{\uline{NB} Si à un carré $Q$ on associe le groupe $\mathfrak{D}=\operatorname{Aut}(Q)$, puis le carré associé à ce groupe, on trouve le carré dont l'ens des sommets est $C\amalg D$, l'action de l'antipodisme dans \deleted{\unsure{chaque}} chacun d'eux $C$ et $D$, i.e. les diagonales sont $C$ et $D$. On trouve donc un carré avec diagonale \deleted{\unsure{privil.}}\inserted{distinguée}\ldots}

Si on part, pour décrire la structure de carré, de l'ens des sommets $S$, cette structure est exprimée par la donnée d'un ens. \inserted{$S$} à $4$ éléments -- i.e. d'un torseur sous un vectoriel $V_S$ de rang $2$ sur $\mathbb{F}_2$ -- et d'une \deleted{\unsure{f}} involution ss pt fixes de $S$, i.e. d'un $\underline{a}_S\in V_S^*$ \deleted{Mais \unsure{ce}}\inserted{(\unsure{i.e.} une} \deleted{\unsure{$V=V_S^*$}} partition de type $(2,2)$ de $S$). Mais pour
% archive foliation: 16
